\appendix
\part{Anhang}

\newcommand{\lit}[3]{\vspace{5pt}\noindent\textbf{[#1]} #2\par\nopagebreak\vspace{1pt}{\small\itshape\color{grau} #3}\par\vspace{2pt}}

\chapter{Formelsammlung und Übersicht}

\section{Begriffe und Formeln}

\begin{merksatz}
\boxtitel{Teil I und II}
\begin{center}
{\small
\renewcommand{\arraystretch}{1.6}
\begin{tabular}{@{}P{4.4cm}P{4.2cm}@{}}
\toprule
$\mathbb{E}[u\mid a] = \sum_\theta u(a,\theta)\Pr[\theta]$ & Erwartungsnutzen (1.1)\\
$s^{*} = \dfrac{c}{(1-r)\,\Delta u}$ & Behandlungsschwelle (1.2)\\
vier Regeln $\Rightarrow$ $u$ existiert & von Neumann (2.1)\\
$x \sim [p:G,\ 1-p:S]$ & Standardlotterie\\
$u$ konkav & Risikoscheu\\
$-a\ln(1-\mathbb{E}[u])$ & Sicherheitsäquivalent\\
$\sum_\theta \Pr[\theta]\max_a u - \max_a \sum_\theta \Pr[\theta]u$ & EVPI (3.1)\\
$\text{EVPI} \ge 0$ & Satz 3.2\\
$p^{*} = \dfrac{C}{C+B}$ & Behandlungsschwelle mit Kosten\\
$p_t$, $p_{tt}$ & Test- und Test-Behandlungs-Schwelle (4.1)\\
Kipppunkt & Empfindlichkeit (5)\\
$\sum_t \sum_z \pi_t(z)\,u(z)\,(1+d)^{-t}$ & QALY im Markow-Modell (6.1)\\
\bottomrule
\end{tabular}}
\end{center}
\end{merksatz}

\begin{merksatz}
\boxtitel{Teil III bis V}
\begin{center}
{\small
\renewcommand{\arraystretch}{1.6}
\begin{tabular}{@{}P{4.4cm}P{4.2cm}@{}}
\toprule
$V_t(z) = \max_a \bigl[r + \mathbb{E}\,V_{t+1}\bigr]$ & Bellman, endlich (7.1)\\
$V = \max_a \bigl[r + \gamma\sum p\,V\bigr]$ & Fixpunkt (8.1)\\
$\gamma$ & Abwertung je Schritt\\
Bedauern $\sim \ln T$ & Bandit, beste Ordnung (9)\\
$\dfrac{\Delta K}{\Delta E}$ & Kosten-Wirksamkeit (10.1)\\
$\lambda\,\Delta E - \Delta K$ & Nettonutzen\\
erweiterte Dominanz & Rangfolge bilden\\
$\max_a \min_\theta u$ & vorsichtige Regel (11)\\
$\max_a \sum_\theta \Pr u$ gegen Intervalle & Ambiguität\\
Bezugspunkt, Verlustscheu & tatsächliches Verhalten (12)\\
keine Regel erfüllt alle Bedingungen & Arrow (13.1)\\
$\bigl(\sum_i q_i^{1-\epsilon}\bigr)^{1/(1-\epsilon)}$ & Ungleichheitsscheu (14)\\
\bottomrule
\end{tabular}}
\end{center}
\end{merksatz}

\section{Die Landkarte}

\begin{merksatz}
\boxtitel{Welcher Satz worauf beruht}
\begin{center}
{\small
\renewcommand{\arraystretch}{1.4}
\begin{tabular}{@{}P{3.6cm}P{1.0cm}P{3.6cm}@{}}
\toprule
\textbf{Ergebnis} & \textbf{Kap.} & \textbf{beruht auf}\\
\midrule
Behandlungsschwelle & $1$ & lineare Nutzenbilanz\\
Nutzenexistenz & $2$ & Stetigkeit und Unabhängigkeit\\
Information schadet nie & $3$ & Maximum ist konvex\\
zwei Testschwellen & $4$ & drei Geraden über $p$\\
Kipppunkt & $5$ & Schnitt zweier Geraden\\
Markow-Bilanz & $6$ & Zustandsverteilung je Zyklus\\
Bellman & $7$ & Zerlegung nach dem ersten Schritt\\
Fixpunkt & $8$ & Kontraktion, Banach (Band 13)\\
Bandit-Schranke & $9$ & Konzentrations\-ungleichung\\
erweiterte Dominanz & $10$ & konvexe Hülle\\
Arrow & $13$ & entscheidende Koalitionen\\
Ungleichheitsmaß & $14$ & Potenzmittel\\
\bottomrule
\end{tabular}}
\end{center}
\end{merksatz}

\section{Was nicht bewiesen wurde}

\begin{merksatz}
\boxtitel{Drei Auslassungen}
\begin{center}
{\small
\renewcommand{\arraystretch}{1.5}
\begin{tabular}{@{}P{3.4cm}P{1.2cm}P{3.4cm}@{}}
\toprule
\textbf{Ergebnis} & \textbf{Kap.} & \textbf{Wo es steht}\\
\midrule
Satz von Arrow & $13$ & Gaertner\zit{9}\\
Bandit-Schranke von Lai und Robbins & $9$ & Lattimore, Szepesvári\zit{16}\\
Konvergenz der Wertiteration & $8$ & Puterman\zit{21}\\
\bottomrule
\end{tabular}}
\end{center}

\begin{center}
\begin{minipage}{0.88\textwidth}
\kursiv{\textbf{Alle übrigen Sätze des Bandes sind vollständig bewiesen} --- darunter Satz 2.1 über die Existenz der Nutzenfunktion, Satz 3.2 über den Wert der Information, Satz 4.1 über die beiden Testschwellen und Satz 7.1 über die Rückwärtsrechnung.}
\end{minipage}
\end{center}
\end{merksatz}

\section{Zeittafel}

\begin{merksatz}
\boxtitel{Von Bernoulli bis zur Nutzenbewertung}
\begin{center}
{\small
\renewcommand{\arraystretch}{1.32}
\begin{tabular}{@{}P{1.3cm}P{3.0cm}P{4.0cm}@{}}
\toprule
\textbf{Jahr} & \textbf{Person} & \textbf{Ereignis}\\
\midrule
$1713$ & N.\ Bernoulli & das Petersburger Spiel\\
$1738$ & D.\ Bernoulli & Nutzen statt Geldwert\\
$1921$ & Knight & Risiko gegen Ungewissheit\\
$1926$ & Ramsey & Wahrscheinlichkeit aus Vorlieben\\
$1944$ & \textbf{von Neumann, Morgenstern} & \textbf{Nutzenaxiome}\\
$1947$ & Wald & Minimax-Regel\\
$1951$ & \textbf{Arrow} & \textbf{Unmöglichkeitssatz}\\
$1953$ & Allais & das Paradox\\
$1954$ & Savage & subjektive Wahrscheinlichkeit\\
$1957$ & Bellman & dynamische Programmierung\\
$1961$ & Ellsberg & Ambiguitätsscheu\\
$1966$ & Howard & Entscheidungsanalyse, Informationswert\\
$1968$ & Klarman u.\,a. & erste QALY-Rechnung\\
$1975$ & Pauker, Kassirer & Behandlungsschwelle\\
$1979$ & Kahneman, Tversky & Neue Erwartungstheorie\\
$1980$ & Pauker, Kassirer & Testschwellen\\
$1985$ & Lai, Robbins & Bandit-Schranke\\
$1990$ & --- & Kosten-Wirksamkeit in der Erstattung\\
$1999$ & NICE & Schwellenwert in der Praxis\\
$2002$ & Kahneman & Nobelpreis\\
$2010$ & --- & probabilistische Modellanalyse als Standard\\
\bottomrule
\end{tabular}}
\end{center}
\end{merksatz}

\section{Glossar}

\begin{merksatz}
\boxtitel{Begriffe von A bis Z}
\begin{center}
{\small
\renewcommand{\arraystretch}{1.3}
\begin{longtable}{@{}P{3.6cm}P{5.0cm}@{}}
\toprule
\textbf{Begriff} & \textbf{Bedeutung}\\
\midrule
\endhead
Ambiguität & unsichere Wahrscheinlichkeiten (11)\\
Bandit & Erkunden gegen Ausnutzen (9.1)\\
Behandlungsschwelle & ab hier behandeln (1.2)\\
Bellman-Gleichung & Wert heute aus Wert morgen (7.1)\\
Bezugspunkt & woran Gewinn gemessen wird (12)\\
Diskontierung & Abwertung künftigen Nutzens (6)\\
Dominanz, erweiterte & Mischung schlägt Einzelne (10)\\
Entscheidungsbaum & Wahl- und Zufallsknoten (5)\\
Erwartungsnutzen & Nutzen mal Wahrscheinlichkeit (1.1)\\
EVPI & Wert vollkommener Information (3.1)\\
EVSI & Wert eines wirklichen Tests (3)\\
ICER & Kosten je zusätzlicher Einheit (10.1)\\
Kipppunkt & wo die Entscheidung wechselt (5)\\
Markow-Modell & Zustände und Zyklen (6.1)\\
Minimax & das schlechteste Ergebnis verbessern (11)\\
Nettonutzen & $\lambda\Delta E - \Delta K$ (10)\\
QALY & Lebensjahr mal Nutzen (6)\\
Risikoprämie & Abstand zum Sicherheitsäquivalent (2)\\
Sicherheitsäquivalent & sicherer Wert gleichen Nutzens (2)\\
Standardlotterie & Nutzenerhebung über $p$ (2)\\
Testschwelle & ab hier testen (4.1)\\
Verlustscheu & Verluste wiegen schwerer (12)\\
Wertiteration & Fixpunkt schrittweise (8.1)\\
Zeitabwägung & Jahre gegen Zustand (2)\\
\bottomrule
\end{longtable}}
\end{center}
\end{merksatz}
